\documentclass[10pt,a4paper]{amsart}
\usepackage{amsmath,amssymb,amsthm,graphicx}
\usepackage{verbatim}
\newenvironment{reference}{\comment}{\endcomment}
%\newenvironment{reference}{}{}
\newenvironment{slogan}{\comment}{\endcomment}
\newenvironment{history}{\comment}{\endcomment}

% For commutative diagrams we use Xy-pic
\usepackage[all]{xy}

% We use 2cell for 2-commutative diagrams.
\xyoption{2cell}
% Theorem environments.
%
\theoremstyle{plain}
\newtheorem{theorem}[subsection]{Theorem}
\newtheorem{proposition}[subsection]{Proposition}
\newtheorem{lemma}[subsection]{Lemma}

\theoremstyle{definition}
\newtheorem{definition}[subsection]{Definition}
\newtheorem{example}[subsection]{Example}
\newtheorem{exercise}[subsection]{Exercise}
\newtheorem{situation}[subsection]{Situation}

\theoremstyle{remark}
\newtheorem{remark}[subsection]{Remark}
\newtheorem{remarks}[subsection]{Remarks}

\numberwithin{equation}{subsection}

% Macros
%
\def\lim{\mathop{\mathrm{lim}}\nolimits}
\def\colim{\mathop{\mathrm{colim}}\nolimits}
\def\Spec{\mathop{\mathrm{Spec}}}
\def\Hom{\mathop{\mathrm{Hom}}\nolimits}
\def\Ext{\mathop{\mathrm{Ext}}\nolimits}
\def\SheafHom{\mathop{\mathcal{H}\!\mathit{om}}\nolimits}
\def\SheafExt{\mathop{\mathcal{E}\!\mathit{xt}}\nolimits}
\def\Sch{\mathit{Sch}}
\def\Mor{\mathop{\mathrm{Mor}}\nolimits}
\def\Ob{\mathop{\mathrm{Ob}}\nolimits}
\def\Sh{\mathop{\mathit{Sh}}\nolimits}
\def\NL{\mathop{N\!L}\nolimits}
\def\CH{\mathop{\mathrm{CH}}\nolimits}
\def\proetale{{pro\text{-}\acute{e}tale}}
\def\etale{{\acute{e}tale}}
\def\QCoh{\mathit{QCoh}}
\def\Ker{\mathop{\mathrm{Ker}}}
\def\Im{\mathop{\mathrm{Im}}}
\def\Coker{\mathop{\mathrm{Coker}}}
\def\Coim{\mathop{\mathrm{Coim}}}

% Boxtimes
%
\DeclareMathSymbol{\boxtimes}{\mathbin}{AMSa}{"02}

%
% Macros for moduli stacks/spaces
%
\def\QCohstack{\mathcal{QC}\!\mathit{oh}}
\def\Cohstack{\mathcal{C}\!\mathit{oh}}
\def\Spacesstack{\mathcal{S}\!\mathit{paces}}
\def\Quotfunctor{\mathrm{Quot}}
\def\Hilbfunctor{\mathrm{Hilb}}
\def\Curvesstack{\mathcal{C}\!\mathit{urves}}
\def\Polarizedstack{\mathcal{P}\!\mathit{olarized}}
\def\Complexesstack{\mathcal{C}\!\mathit{omplexes}}
% \Pic is the operator that assigns to X its picard group, usage \Pic(X)
% \Picardstack_{X/B} denotes the Picard stack of X over B
% \Picardfunctor_{X/B} denotes the Picard functor of X over B
\def\Pic{\mathop{\mathrm{Pic}}\nolimits}
\def\Picardstack{\mathcal{P}\!\mathit{ic}}
\def\Picardfunctor{\mathrm{Pic}}
\def\Deformationcategory{\mathcal{D}\!\mathit{ef}}
\usepackage{cleveref}
\crefname{originalremark}{original remark}{original remarks}
\Crefname{originalremark}{Original remark}{Original remarks}
\makeatletter
\newlabel{remark-morphism-topoi-comes-from-morphism-sites}{{03CJ}{source}{}{sourcepointer-1}{}}
\newlabel{remark-morphism-topoi-comes-from-morphism-sites@cref}{{[originalremark][0][]03CJ}{[1][0][]source}{}{}{}}
\makeatother
\crefname{originalproposition}{original proposition}{original propositions}
\Crefname{originalproposition}{Original proposition}{Original propositions}
\makeatletter
\newlabel{proposition-get-morphism}{{00X6}{source}{}{sourcepointer-2}{}}
\newlabel{proposition-get-morphism@cref}{{[originalproposition][0][]00X6}{[1][0][]source}{}{}{}}
\makeatother
\crefname{originallemma}{original lemma}{original lemmas}
\Crefname{originallemma}{Original lemma}{Original lemmas}
\makeatletter
\newlabel{lemma-pullback-representable-sheaf}{{04D3}{source}{}{sourcepointer-3}{}}
\newlabel{lemma-pullback-representable-sheaf@cref}{{[originallemma][0][]04D3}{[1][0][]source}{}{}{}}
\makeatother
\crefname{originallemma}{original lemma}{original lemmas}
\Crefname{originallemma}{Original lemma}{Original lemmas}
\makeatletter
\newlabel{lemma-localize-morphism}{{03CF}{source}{}{sourcepointer-4}{}}
\newlabel{lemma-localize-morphism@cref}{{[originallemma][0][]03CF}{[1][0][]source}{}{}{}}
\makeatother
\crefname{originallemma}{original lemma}{original lemmas}
\Crefname{originallemma}{Original lemma}{Original lemmas}
\makeatletter
\newlabel{lemma-localize-compare}{{04IQ}{source}{}{sourcepointer-5}{}}
\newlabel{lemma-localize-compare@cref}{{[originallemma][0][]04IQ}{[1][0][]source}{}{}{}}
\makeatother
\title{Localizing a morphism of topoi}
\author{The Stacks Project Authors}
\date{stacks-project:sites.tex:local-snapshot-20261012; GNU FDL 1.2 or later}
\begin{document}
\maketitle
\noindent\textit{Editorial note.} All delivered mathematical text is the author's own, copied from the pinned Stacks Project chapter file sites-cohomology.tex; no statement or proof step was rewritten, reordered or removed.\par\medskip
\noindent\textit{Editorial note.} The proof is the author's complete argument; every step the author writes out is inside the counted span and no closing sentence is a placeholder.\par\medskip
\noindent\textit{Editorial note.} Every labelled reference occurring inside the counted spans is carried in assembly.external\_references as a non-load-bearing navigation pointer (same-chapter ones at their definition site, cross-chapter ones at the invocation inside the counted span), so no load-bearing mathematics is replaced by a stub. No bibliography citation occurs inside the counted spans.\par\medskip
\noindent\textit{Editorial note.} The checked chapter list for same-name collisions is the set of stacks chapters staged in this repository: chow, cohomology, divisors, etale-cohomology, flat, more-algebra, more-morphisms, morphisms, spaces-more-morphisms, varieties together with this card's chapter sites-cohomology.\par\medskip
\noindent\textit{Editorial note.} The preamble is the chapter's own retained minimal preamble (sources/stacks-sites-cohomology/retained-preamble.tex), a verbatim excerpt of the native preamble: the Xy-pic setup and the native theorem declarations with the mathematical macros.\par\medskip
\noindent\textit{Allowed standard prerequisites (editorial).}\begin{itemize}
\item Replacing topoi by sites. Same chapter, Remark remark-morphism-topoi-comes-from-morphism-sites. Allows the morphism to be assumed to come from a continuous functor between suitable sites.
\item The proposition producing a morphism of sites. Same chapter, Proposition proposition-get-morphism. Supplies the assumptions on the functor.
\item Localizing a morphism of sites. Same chapter, Lemma lemma-localize-morphism. Supplies the commutative diagram between the localized sites.
\item The pullback of a representable sheaf. Same chapter, Lemma lemma-pullback-representable-sheaf. Identifies the pullback of the localizing sheaf with a representable one.
\end{itemize}\medskip
\ifdefined\hypertarget\hypertarget{sourcepointer-1}{}\fi
\noindent\textit{Source reference.} Original remark 03CJ: sources/stacks-sites/pinned-original-source.tex, source lines 7868--7868. Navigation pointer invoked at proof line 7868; the counted span contains the pointer only and the cited result is not reproduced inside it, so the pointer is registered as a non-load-bearing navigation stub in the established form for Stacks material. No bibliography citation is involved.\par\medskip
\ifdefined\hypertarget\hypertarget{sourcepointer-2}{}\fi
\noindent\textit{Source reference.} Original proposition 00X6: sources/stacks-sites/pinned-original-source.tex, source lines 7871--7871. Navigation pointer invoked at proof line 7871; the counted span contains the pointer only and the cited result is not reproduced inside it, so the pointer is registered as a non-load-bearing navigation stub in the established form for Stacks material. No bibliography citation is involved.\par\medskip
\ifdefined\hypertarget\hypertarget{sourcepointer-3}{}\fi
\noindent\textit{Source reference.} Original lemma 04D3: sources/stacks-sites/pinned-original-source.tex, source lines 7875--7875. Navigation pointer invoked at proof line 7875; the counted span contains the pointer only and the cited result is not reproduced inside it, so the pointer is registered as a non-load-bearing navigation stub in the established form for Stacks material. No bibliography citation is involved.\par\medskip
\ifdefined\hypertarget\hypertarget{sourcepointer-4}{}\fi
\noindent\textit{Source reference.} Original lemma 03CF: sources/stacks-sites/pinned-original-source.tex, source lines 7877--7877. Navigation pointer invoked at proof line 7877; the counted span contains the pointer only and the cited result is not reproduced inside it, so the pointer is registered as a non-load-bearing navigation stub in the established form for Stacks material. No bibliography citation is involved.\par\medskip
\ifdefined\hypertarget\hypertarget{sourcepointer-5}{}\fi
\noindent\textit{Source reference.} Original lemma 04IQ: sources/stacks-sites/pinned-original-source.tex, source lines 7888--7888. Navigation pointer invoked at proof line 7888; the counted span contains the pointer only and the cited result is not reproduced inside it, so the pointer is registered as a non-load-bearing navigation stub in the established form for Stacks material. No bibliography citation is involved.\par\medskip
\begin{lemma}
\label{lemma-localize-morphism-topoi}
Let $f : \Sh(\mathcal{C}) \to \Sh(\mathcal{D})$
be a morphism of topoi. Let $\mathcal{G}$ be a sheaf on $\mathcal{D}$.
Set $\mathcal{F} = f^{-1}\mathcal{G}$. Then there exists
a commutative diagram of topoi
$$
\xymatrix{
\Sh(\mathcal{C})/\mathcal{F} \ar[r]_{j_\mathcal{F}} \ar[d]_{f'} &
\Sh(\mathcal{C}) \ar[d]^f \\
\Sh(\mathcal{D})/\mathcal{G} \ar[r]^{j_\mathcal{G}} &
\Sh(\mathcal{D}).
}
$$
The morphism $f'$ is characterized by the property that
$$
(f')^{-1}(\mathcal{H} \xrightarrow{\varphi} \mathcal{G})
=
(f^{-1}\mathcal{H} \xrightarrow{f^{-1}\varphi} \mathcal{F})
$$
and we have $f'_*j_\mathcal{F}^{-1} = j_\mathcal{G}^{-1}f_*$.
\end{lemma}

\begin{proof}
Since the statement is about topoi and does not refer to the underlying
sites we may change sites at will. Hence by the discussion in
Remark \ref{remark-morphism-topoi-comes-from-morphism-sites}
we may assume that $f$ is given by a continuous functor
$u : \mathcal{D} \to \mathcal{C}$ satisfying the assumptions of
Proposition \ref{proposition-get-morphism}
between sites having all finite limits and subcanonical topologies,
and such that $\mathcal{G} = h_V$ for some object $V$ of
$\mathcal{D}$. Then $\mathcal{F} = f^{-1}h_V = h_{u(V)}$ by
Lemma \ref{lemma-pullback-representable-sheaf}.
By
Lemma \ref{lemma-localize-morphism}
we obtain a commutative diagram of morphisms of topoi
$$
\xymatrix{
\Sh(\mathcal{C}/U) \ar[r]_{j_U} \ar[d]_{f'} &
\Sh(\mathcal{C}) \ar[d]^f \\
\Sh(\mathcal{D}/V) \ar[r]^{j_V} &
\Sh(\mathcal{D}),
}
$$
and we have $f'_*j_U^{-1} = j_V^{-1}f_*$. By
Lemma \ref{lemma-localize-compare}
we may identify $j_\mathcal{F}$ and $j_U$
and $j_\mathcal{G}$ and $j_V$. The description of $(f')^{-1}$
is given in
Lemma \ref{lemma-localize-morphism}.
\end{proof}

\end{document}
